\properlane{song-exile}{Learning thanksgiving in exile}

Ephesians commands thanksgiving and song; the Offertory supplies tears.
The question is not which voice the assembly should silence. It is how
both can become faithful speech. Gratitude that cannot acknowledge
suffering would fail the captives beside Babylon's rivers. Lament that
could no longer hope would fail the servant comforted by God's word.
The Mass forms a community that can confess its sin, grieve its losses,
receive its food, and praise its Lord without confusing those acts.

\subsection{Evil days do not make creation evil}
Chrysostom explains the Epistle's evil days by separating created time
from the wrongdoing committed within it. Day is God's creature; evil
belongs to conduct. The distinction prevents the injunction to redeem
time from becoming despair about the world itself. It also makes the
command active. Circumstances may be hostile, but believers still have
to discern how to walk within them.\footnote{Chrysostom, \work{Homilies
on Ephesians}, XIX, exposition of 5:15--17.}

The care he urges is neither cowardice nor a search for quarrels.
Christians should not add needless provocations to the offence occasioned
by the Gospel. They can surrender lesser things without surrendering
fidelity. This interpretation gives patience an intelligent object:
preserving the good that matters in circumstances where anger would
otherwise decide the next action. Redeeming time means making an
opportunity serve that good, even when one cannot make the whole
situation favorable.

This community begins the Mass with Azarias's confession. Jerome's
representative reading matters particularly here. The youths can
acknowledge their people's sins while suffering for their own refusal
of idolatry. They belong to the people whose need they confess; their
faithfulness does not establish a private refuge from solidarity.
\footnote{Jerome, \work{Commentary on Daniel}, I, at 3:29.}
At the same time, solidarity does not erase distinctions of responsibility.
The king's violence remains the king's violence. Confession in the first
person plural does not authorize accusing each injured member of the
community of causing the injury.

The Collect asks what such a people needs in order to serve together:
pardon and peace. A congregation bound by resentment would reproduce
within itself the conduct it deplores around it. Schuster's heart divided
by its passions is not only an individual difficulty. Many divided
hearts can make common action impossible. Cleansing and peaceful
service therefore concern the shared life as well as the inward life.
The Introit's blessed way and the Epistle's wise walking both lead
beyond a moment of sentiment into a sustained manner of living.

\subsection{Thanks amid affliction: the strength and difficulty of Chrysostom's answer}
Chrysostom takes continual thanksgiving very far. It includes sickness,
poverty, anguish, and blessings not yet recognized as blessings. Giving
thanks only in prosperity would reveal little about trust; the test
arrives when visible benefit cannot settle the question. He appeals to
God's goodness toward people who do not even recognize their giver and
to the greater love of the heavenly Father.\footnote{Chrysostom,
\work{Homilies on Ephesians}, XIX, exposition of 5:20.}
His argument rests on confidence in divine providence, not on an assertion
that affliction feels agreeable.

There is a real difference between this reading and the nineteenth-century
editorial note in the NPNF translation. Note~404 restricts the expression
``all things'' to blessings, against Chrysostom's extension. The note is
an editor's judgment, not another sentence of the homily. Chrysostom's
exposition remains the more demanding one: gratitude must survive what
the petitioner neither chose nor presently understands.
\footnote{NPNF I.13, \work{Homilies on Ephesians}, XIX, editorial note~404.}
That demand must be heard alongside his distinction between evil conduct
and created days. Thanksgiving to God amid affliction does not transform
wicked action into a virtue or make the wrongdoer worthy of thanks for it.

The wider formulary sharpens that distinction. The Introit confesses
disobedience as sin; the Secret asks that vices be expelled. A thanksgiving
that approved vice would contradict the prayers with which the Church
gives thanks. The Gospel's father actively seeks help for an endangered
child. He does not demonstrate faith by allowing a preventable harm to
continue. The Offertory laments a devastated home. Its tears remain a
truthful part of the offering.

Chrysostom's own rhetoric sometimes minimizes bodily afflictions in order
to distinguish them from sin. The Gospel's urgency gives bodily life its
full weight here: the son is dying, and the petition concerns his life.
The moral distinction need not become a denial of pain. The assembly's
thanksgiving can confess God's goodness, seek relief, repent of wrongdoing,
and accompany sufferers without claiming to know the hidden reason for
each affliction. Its first duty is fidelity, not an explanation imposed
on another person's wound.

\subsection{Beside the currents, with Zion remembered}
Augustine asks what the rivers of Babylon signify for the pilgrim Church.
His answer concerns passing loves that can carry a person away: wealth,
status, and the pleasures whose movement promises stability but cannot
give it. The faithful sit beside the rivers. Their position becomes an
image of humility and resistance. They recognize where they are without
consenting to be swept along.\footnote{Augustine, \work{Expositions of
the Psalms}, Ps.~136, \S\S2--4 (NPNF heading Ps.~137).}

He then distinguishes two sorts of sorrow. One can weep because worldly
gain has been lost and still desire only the same gain. Or one can weep
in remembrance of Zion, measuring present life by an abiding good.
The latter grief can exist even when life in Babylon appears successful.
The distinction is not between people with feelings and people without
them. It concerns what their feelings reveal about the home they seek.
Memory becomes a discipline of love.

The historical lament remains the prayer of captives remembering
Jerusalem. Augustine's Christian exposition addresses those whose
captivity is also moral and spiritual. In the full psalm the captors'
demand for music turns singing into a question of allegiance. Not every
occasion for a song is an occasion for truthful praise. Refusing a
performance for mockery is different from forgetting the song itself.
That difference allows Augustine to finish by urging the pilgrims to
sing the songs of Zion to one another.
\footnote{Augustine, Ps.~136, \S\S5--9,13.}

His conclusion meets Ephesians at a precise point: song passes between
members of the community and sustains their movement toward God.
Chrysostom's inward attention prevents the song from becoming an empty
exercise; Augustine's remembrance prevents it from becoming accommodation
to whatever the world presently rewards. The Alleluia's prepared heart
then means more than musical readiness. It is a heart whose praise has
an object firm enough to survive changing circumstances.

Bellarmine develops the distinction between location and allegiance.
A person may hold responsibilities in the world without giving the
heart over to its passing goods. His exposition can therefore include
those in positions of rule who regard their authority as a burden of
service rather than a possession to admire.\footnote{Bellarmine,
\work{Commentary on the Psalms}, Ps.~136:1--4.}
Longing for Zion does not require neglecting one's household or abandoning
the duties through which others depend on one's care. It changes how
those duties are performed and what reward is sought from them.

\subsection{The open hand makes praise responsible}
The Gradual's universal language places a demanding claim between the
Epistle's instruction and the Gospel's particular cure. All creatures
look toward the divine giver. Yet some creatures hunger. Augustine
answers the question of delay through the fitting time known by the
physician, and the final salvation toward which he directs the patient.
Bellarmine adds a moral question addressed to the person who has enough:
what has human hoarding done to the distribution of God's gifts?
\footnote{Augustine, Ps.~144, \S\S14--17; Bellarmine, Ps.~144:15--16.}

The two questions have different addressees. The sufferer is encouraged
not to surrender hope; the possessor is challenged not to withhold help.
A comfortable person cannot use the first answer to evade the second.
To praise the open divine hand while closing one's own would make praise
contradict conduct. Bellarmine's account gives generosity a theological
shape: human care should resemble the giver on whom everyone depends.
His further conjectures about a sufferer's fault do not identify the
cause of any particular person's hunger. The demand to share remains
clear without such a diagnosis.

Chrysostom's exposition of mutual submission gives this sharing a
communal form. The master who imagines himself only a recipient of
service must consider the maintenance he owes those who work. The
ideal he presses is voluntary interchange, illustrated by friends
serving one another. The inward song of verse~19 and the reciprocal
conduct of verse~21 cannot be separated without cutting the Epistle
short. In thanksgiving the assembly acknowledges a giver; in mutual
service it lets that acknowledgment change its common life.

The Gospel's ending extends the horizon from one petition to a whole
household's belief. The father comes on behalf of another, receives
the servants' testimony, and shares in a faith that includes the house.
Gregory's comparison with the centurion rebukes deference based on
prestige. Christ does not value a ruler's child because rank makes that
child more human than a servant.\footnote{Gregory, II.28.2; Schuster,
III, pp.~176--177, on the household's faith and its apostolic consequence.}
The community formed by the healing word has reason to listen to
servants, care for dependents, and refuse the honor system that sees
only the possessions surrounding a person.

\subsection{A healed community receives a greater home}
The Secret turns the community's attention inward without privatizing
it. Vices of the heart have effects beyond the heart: pride resists
service, avarice withholds food, resentment makes common song hollow.
The petition for heavenly medicine asks Christ to heal the sources of
these fractures. Schuster's Eucharistic account keeps the remedy's
agent and depth clear. The sacrament is God's gift for the whole person,
not simply a technique for producing agreeable social behavior.

The Communion returns to the word that gives hope. Augustine applies
its consolation to Christ's body under persecution and carries its
promise toward resurrection. A humiliated people can ask for fulfilment
because the promise precedes their success. His following exposition
speaks of God's statutes as songs in the house of pilgrimage. Even
where life has the character of night, the Lord's name can be remembered.\footnote{Augustine, Ps.~118, \S\S50--56. The songs of pilgrimage are
the context of verse~54, not additional words of the Communion.}
This remembered promise supplies continuity between the song commanded
in the Epistle and the home remembered in the Offertory.

The Postcommunion asks that this continuity become obedience. A
community's thanksgiving has to endure after the act of singing ends.
Its members need the grace to keep commandments in their dealings with
one another. Since the prayer asks God to grant that obedience, its
vision of a renewed common life is a response to grace, not a condition
that the community must first manufacture in order to purchase grace.

The commemorative Marian prayers, when appointed, locate this hope in
the Word's actual human birth, petition for present and perpetual peace,
and ask cleansing and fellowship in heavenly remedy. They strengthen
the pilgrim people's dependence on Christ while acknowledging Mary's
intercession. The Trinity Preface directs the assembly beyond its own
social unity toward the God it worships: the Father, Son, and Spirit
receive one adoration in equal majesty. The assembly's unity is a gift
received in worship, not an alternative object of worship.

Augustine's Zion is finally the place where brotherhood is secure,
temptation no longer divides, and death is conquered. Schuster's
Eucharistic nourishment likewise tends toward vision. That final home
gives present generosity its hope: feeding a neighbor is neither the
whole of salvation nor an irrelevant task beside it. It belongs to
the earthly obedience of those who desire to live together with God.
The song can be thankful now because its giver is faithful, and it
can include tears now because its final fullness has not yet arrived.

\begin{fourSenses}
\item[Literal.]
Azarias's prayer and the Babylonian lament concern a people in trial;
Ephesians addresses a community called to wise conduct, song, and mutual
submission. John describes a particular household's healing and belief.
Thanksgiving and lament both stand among the appointed words.
\item[Allegorical.]
Augustine's pilgrims remember Zion as their promised home, while Christ
forms a believing household and heals his Church. The incarnate Word,
confessed also in the commemorative branch, nourishes this people on its
journey to the Father in the Holy Spirit.
\item[Moral.]
Evil conduct is resisted through prudent fidelity, repentance,
attentive praise, reciprocal service, and generosity. Bellarmine's open
hand tests the praise of those who possess. Lament tells the truth about
loss; gratitude to God does not require approval of cruelty or vice.
\item[Anagogical.]
Zion's lasting fellowship, resurrection, and the vision of God complete
the pilgrim song. Present gifts sustain hope for that common home.
Responsibility for others belongs to the journey; hope for its end
does not dispense with service along the way.
\end{fourSenses}
